\documentclass[11pt,a4paper]{article}

\usepackage[utf8]{inputenc}
\usepackage[T1]{fontenc}
\usepackage{amsmath,amssymb,amsthm}
\usepackage{booktabs}
\usepackage{longtable}
\usepackage{hyperref}
\usepackage[margin=1in]{geometry}
\usepackage{natbib}
\usepackage{xcolor}
\usepackage{enumitem}
\usepackage{caption}
\usepackage{array}
\usepackage{tabularx}
\usepackage{multirow}
\usepackage{tikz}
\usetikzlibrary{arrows.meta, positioning, calc, fit, backgrounds,
                shapes.geometric, decorations.pathreplacing, decorations.pathmorphing}
% All figures in this draft use manual TikZ positioning, so we stay on pdfLaTeX.
% To switch to auto-laid-out figures (Fig. 2, Fig. 4), uncomment the two lines
% below AND compile with LuaLaTeX (lualatex main.tex). See figure-source comments.
% \usetikzlibrary{graphs,graphdrawing}
% \usegdlibrary{layered,force}
\usepackage{pgfplots}
\pgfplotsset{compat=1.18}

% notation.tex — single source of truth for active-inference macros.
% % notation.tex — single source of truth for active-inference macros.
% % notation.tex — single source of truth for active-inference macros.
% \input{notation} from main.tex preamble. Subagents must use only these macros
% and must not edit this file.

\newcommand{\FE}{\mathcal{F}}                    % variational free energy
\newcommand{\EFE}{\mathcal{G}}                   % expected free energy
\newcommand{\KL}[2]{D_{\mathrm{KL}}\!\left[#1 \,\|\, #2\right]}
\newcommand{\E}[1]{\mathbb{E}\!\left[#1\right]}
\newcommand{\Eunder}[2]{\mathbb{E}_{#1}\!\left[#2\right]}

% Active inference primitives
\newcommand{\C}{\mathbf{C}}                      % prior preferences (log-preferences)
\newcommand{\Prec}{\boldsymbol{\Pi}}             % precision matrix / vector
\newcommand{\prc}{\gamma}                        % scalar precision (avoid \prec — already taken)
\newcommand{\policy}{\pi}                        % policy
\newcommand{\hidden}{s}                          % hidden state
\newcommand{\obs}{o}                             % observation
\newcommand{\Q}{Q}                               % approximate posterior
\newcommand{\Pgen}{P}                            % generative model

% Domain-specific
\newcommand{\salience}{\mathrm{Sal}}             % salience operator
\newcommand{\CRED}{\mathrm{CRED}}                % credibility-enhancing display
\newcommand{\norm}{\mathcal{N}}                  % a norm (region of high-precision shared C)

% Drafting
\newcommand{\todo}[1]{\textcolor{red}{\textbf{[TODO: #1]}}}
 from main.tex preamble. Subagents must use only these macros
% and must not edit this file.

\newcommand{\FE}{\mathcal{F}}                    % variational free energy
\newcommand{\EFE}{\mathcal{G}}                   % expected free energy
\newcommand{\KL}[2]{D_{\mathrm{KL}}\!\left[#1 \,\|\, #2\right]}
\newcommand{\E}[1]{\mathbb{E}\!\left[#1\right]}
\newcommand{\Eunder}[2]{\mathbb{E}_{#1}\!\left[#2\right]}

% Active inference primitives
\newcommand{\C}{\mathbf{C}}                      % prior preferences (log-preferences)
\newcommand{\Prec}{\boldsymbol{\Pi}}             % precision matrix / vector
\newcommand{\prc}{\gamma}                        % scalar precision (avoid \prec — already taken)
\newcommand{\policy}{\pi}                        % policy
\newcommand{\hidden}{s}                          % hidden state
\newcommand{\obs}{o}                             % observation
\newcommand{\Q}{Q}                               % approximate posterior
\newcommand{\Pgen}{P}                            % generative model

% Domain-specific
\newcommand{\salience}{\mathrm{Sal}}             % salience operator
\newcommand{\CRED}{\mathrm{CRED}}                % credibility-enhancing display
\newcommand{\norm}{\mathcal{N}}                  % a norm (region of high-precision shared C)

% Drafting
\newcommand{\todo}[1]{\textcolor{red}{\textbf{[TODO: #1]}}}
 from main.tex preamble. Subagents must use only these macros
% and must not edit this file.

\newcommand{\FE}{\mathcal{F}}                    % variational free energy
\newcommand{\EFE}{\mathcal{G}}                   % expected free energy
\newcommand{\KL}[2]{D_{\mathrm{KL}}\!\left[#1 \,\|\, #2\right]}
\newcommand{\E}[1]{\mathbb{E}\!\left[#1\right]}
\newcommand{\Eunder}[2]{\mathbb{E}_{#1}\!\left[#2\right]}

% Active inference primitives
\newcommand{\C}{\mathbf{C}}                      % prior preferences (log-preferences)
\newcommand{\Prec}{\boldsymbol{\Pi}}             % precision matrix / vector
\newcommand{\prc}{\gamma}                        % scalar precision (avoid \prec — already taken)
\newcommand{\policy}{\pi}                        % policy
\newcommand{\hidden}{s}                          % hidden state
\newcommand{\obs}{o}                             % observation
\newcommand{\Q}{Q}                               % approximate posterior
\newcommand{\Pgen}{P}                            % generative model

% Domain-specific
\newcommand{\salience}{\mathrm{Sal}}             % salience operator
\newcommand{\CRED}{\mathrm{CRED}}                % credibility-enhancing display
\newcommand{\norm}{\mathcal{N}}                  % a norm (region of high-precision shared C)

% Drafting
\newcommand{\todo}[1]{\textcolor{red}{\textbf{[TODO: #1]}}}


% Give figures more breathing room
\setlength{\textfloatsep}{20pt plus 4pt minus 4pt}
\setlength{\intextsep}{18pt plus 4pt minus 4pt}
\setlength{\floatsep}{18pt plus 4pt minus 4pt}
\setlength{\abovecaptionskip}{10pt}
\setlength{\belowcaptionskip}{6pt}

\newtheorem{definition}{Definition}[section]
\newtheorem{proposition}[definition]{Proposition}

\title{\textbf{Norms as Shared Precision Priors: \\[0.3em]
Values, Salience, and Cultural Transmission \\[0.3em]
in Active Inference}}
\author{Jonas Hallgren\textsuperscript{1} \\
\textsuperscript{1}Equilibria Network \\
\texttt{contact@eq-network.org}}
\date{\today \\ \textit{Working Draft --- v0.1}}

\begin{document}

\maketitle

\begin{abstract}
Norms --- shared regularities of behaviour that enable coordination among
agents with heterogeneous preferences --- are studied separately across game
theory, cultural evolution, cognitive sociology, and active inference, with
limited integration. The substrate problem is how prior preferences come to be
shared across a community in the first place: how values are installed in
individual agents through community-level dynamics. We address this problem
within the active-inference framework. Values are identified with prior
preferences ($\C$) in an agent's generative model; prestige cues and
credibility-enhancing displays are formalised as precision modulators on the
channel through which $\C$ is updated from observed behaviour; and norms are
characterised as stabilised regions of high-precision shared $\C$ across a
community. The resulting fixed-point dynamics --- community priors shape
salience, salience directs attention and copying, copying installs priors in
next-generation agents --- admit three regimes (convergence on a shared norm,
bifurcation into multiple sub-community attractors, and dissolution into
incoherence) determined by the precision gain of the social channel. The
shape of $\C$ itself further governs the exploration--exploitation tradeoff
in installed values. The account is complementary to existing treatments of
social influence as policy-precision modulation, operating one layer up on
the content of prior preferences. We propose a multi-agent simulation that
operationalises the regime predictions and discuss implications for tracking
how mediated information environments may restructure community-level value
distributions.
\end{abstract}

%% ============================================================
\section{Introduction}
\label{sec:intro}
%% ============================================================

The cognitive and cultural substrates of norms --- shared regularities of
behaviour through which agents with heterogeneous preferences coordinate
---  have been the subject of sustained inquiry across multiple disciplines,
with limited integration. Game theory and the study of conventions, descended
from Lewis, frame norms as equilibria of repeated interaction
\citep{koster2022SillyRules,oldenburgZhiXuan2024BayesianRule,
vinitsky2023sanctions}. Cultural evolution
\citep{henrichGilWhite2001Prestige,heyes2018CognitiveGadgets} documents the
population-level mechanisms by which behaviours and values are transmitted
across generations, with particular attention to prestige-biased copying and
credibility-enhancing displays \citep{henrich2009CREDs,chudek2012BystanderGaze}.
Cognitive sociology, anchored in Bourdieu's account of habitus and developed
empirically by \citet{vaisey2009DualProcess}, locates norms in durable
dispositions acquired through participation in social life. Active inference
\citep{friston2015ActiveInferenceEpistemic,parrFriston2017Attention} provides
a computational vocabulary in which values are prior preferences in an
agent's generative model and attention is precision-weighting on prediction
error; multi-agent extensions of this framework
\citep{albarracin2022EpistemicCommunities,heins2024CollectiveSurprise} have
begun to apply these tools to collective phenomena.

What has been missing is a formal account of the substrate that connects
these traditions: the cognitive mechanism by which the prior preferences
that make a norm binding for an individual agent come to be shared across a
community. Equilibrium framings yield stable behaviour without specifying a
cognitive substrate; cultural evolution describes the empirical transmission
process in informal vocabulary; cognitive sociology characterises the
phenomenology without committing to a computational mechanism; active
inference has the formal apparatus but has not been systematically applied
at the junction where individual generative models are coupled by social
observation. The substrate problem is the gap into which this paper places
its contribution.

Working in the active-inference framework, we identify values with prior
preferences ($\C$) in an agent's generative model. We propose that prestige
cues and credibility-enhancing displays modulate the precision of the
channel through which $\C$ is updated from observed behaviour, and that
norms are best characterised as stabilised regions of high-precision shared
$\C$ across a community. The resulting fixed-point loop --- community priors
shape salience, salience directs attention and copying, copying installs
priors in next-generation agents --- admits three regimes (convergence on a
shared norm, bifurcation into echo chambers, dissolution into incoherence)
determined by the precision gain of the social channel. The shape of $\C$
itself --- sharp, broad, or what \citet{torresan2025ShapeC} call ``shaped''
--- governs the exploration--exploitation tradeoff and the robustness of
installed values.

The contribution is threefold. First, we provide a formal definition of
norms as high-precision shared regions of $\C$, with explicit stability
conditions for the loop dynamics. Second, we differentiate the account from
\citet{albarracin2025SocialPower}'s treatment of social power as
policy-precision modulation: we operate one layer up, on the content of
prior preferences. Third, we propose a minimum viable multi-agent simulation
that operationalises the regime predictions and provides a target for
empirical work.

The paper proceeds as follows. Section~\ref{sec:cooperative-ai} situates the
substrate problem within contemporary norm research in multi-agent systems
and cognitive science. Section~\ref{sec:cultural-evolution} surveys the
cultural-evolution backdrop, paying particular attention to the
domain-generality of prestige cues, which behave as precision parameters
rather than content priors. Section~\ref{sec:formal} develops the formal
account, including the definition of norms and the bifurcation conditions
of the loop. Section~\ref{sec:simulation} proposes a multi-agent simulation
that operationalises the framework's regime predictions.
Section~\ref{sec:discussion} closes with limitations and remarks on the
implications for environments in which artificial agents mediate the
substrate of value transmission.

\section{Norm Learning Without a Substrate: The Installation Problem}\label{sec:cooperative-ai}

%% ~1500 words / ~2 pages. No EFE formalism. No §3 apparatus.
%% Anchor keys: dafoe2020CooperativeAI, critchKrueger2020ARCHES,
%%   koster2022SillyRules, oldenburgZhiXuan2024BayesianRule,
%%   kulveit2025GradualDisempowerment, albarracin2025SocialPower.
%% Light forward refs: albarracin2022EpistemicCommunities, henrich2009CREDs.

Cooperative intelligence --- the capacity of agents to find mutually
beneficial joint courses of action --- is the organising concept behind a
rapidly consolidating research field. \citet{dafoe2020CooperativeAI} lay out
the agenda: cooperation rests on understanding, communication, commitment, and
norms. Of these four pillars, norms carry the heaviest load. They stabilise
expectations across agents who cannot fully observe one another's preferences,
sustain coordination under incomplete contracting, and provide the decentralised
scaffolding that makes institutions tractable. Whether the agents in question are
AI--AI, AI--human, or human--human-via-AI, a world with more capable AI systems
is a world in which norms do more work, not less~\citep{critchKrueger2020ARCHES}.
Getting their computational basis right is therefore not a peripheral concern for
cooperative AI: it is foundational.

\subsection*{Norms-as-Equilibria: What Works and What Is Missing}

The dominant treatment of norms in cooperative AI follows the game-theoretic
tradition inaugurated by Lewis and formalised by Schelling: a norm is a
self-enforcing equilibrium of behaviour sustained by mutual expectations. In this
frame, a social norm \textit{exists} when compliance and sanctioning behaviour
jointly persist at population scale under coordination or social-dilemma
pressure. The practical advantage of this framing is substantial. It interfaces
naturally with multi-agent reinforcement learning (MARL), makes the stability
conditions of norms tractable, and generates testable predictions about which
equilibria persist under evolutionary pressure. The Leibo lineage of cooperative
MARL papers, the decentralised sanctioning work of \citet{vinitsky2023sanctions},
and the Bayesian rule-induction approach of \citet{oldenburgZhiXuan2024BayesianRule}
all inherit this framing, and all deliver real scientific progress within it.

What the equilibrium framing does not deliver is a \textit{cognitive substrate}
--- a mechanistic story about what a norm is in the head of a single agent.
The equilibrium frame treats the agent's internal representation as a black box:
whatever it contains, if compliance and sanctioning persist, the norm is present.
Three observations mark the limits of this silence. First, violations of norms
\textit{feel surprising} in a way that violations of mere regularities do not;
they elicit moral outrage, not just strategic recalculation. Second, some norms
are \textit{sticky}: agents continue to follow them even when enforcement is
relaxed and even when they cannot articulate a welfare rationale. Third, and
most relevant here, agents can shift from \textit{imitating} a norm --- following
it to avoid sanctions --- to \textit{embodying} it, at which point the norm shapes
what they attend to, what they notice as salient, and what they find imaginable
as a course of action. The equilibrium account has no vocabulary for this
phenomenology, because it operates at the level of policy outputs, not generative
model structure.

The \citet{koster2022SillyRules} result makes this gap vivid in a particularly
elegant way. Adding arbitrary ``silly rules'' --- taboos that carry no welfare
justification --- into a MARL foraging environment \textit{improves} agents'
acquisition of compliance and enforcement with welfare-relevant taboos. This
result is surprising under the equilibrium frame, because silly rules add no
information about the welfare-relevant equilibrium; under a precision-substrate
account, it is exactly what one would expect. Silly rules supply additional
Bayesian evidence that a \textit{normative regime is in operation} ---
they pump up the precision of the agents' prior over the policy distribution
``follow taboos'' without specifying any particular content. Once that
precision is elevated, the agents learn welfare-relevant norms faster. The
mechanism Koster et al.\ identify is not primarily about content; it is about
the legibility of the normative substrate as such. This is the empirical
signature of a precision-shaping story, written in MARL language.

\citet{oldenburgZhiXuan2024BayesianRule} is the most cognitively rich current
treatment and the most direct dialogue partner for this paper. Their agents
Bayesianly induce shared norms as discrete obligative or prohibitive rules from
observed compliance and sanction data; they show that a \textit{shared
assumption} that a norm system is operating is sufficient to bootstrap common
knowledge of specific norms, enabling rapid convergence and stable transmission
to newcomers. This is a genuine advance: unlike equilibrium-only accounts, it
gives norms propositional structure in individual agents. But the ``shared
assumption'' itself remains unmodelled prior content. It is posited as an
external parameter of the inference problem rather than as something agents
acquire from a community, transmit with fidelity, and maintain through
coordinated attention. Representing norms as discrete inferred rules also
forecloses the graded, phenomenologically thick account of norm-internationalisation
that the sticky-norm and violation-surprise phenomena require. The rules are
the outputs of inference; the question of what \textit{installs the prior} that
makes the shared assumption credible is left open.

\subsection*{The Installation Question}

Where do an agent's priors over what-counts-as-a-norm come from? This is the
question the equilibrium and Bayesian-rule accounts both treat as exogenous.
Norms are handed to agents, or emerge from iterated play, but in neither case
does the account say how one generation's normative landscape gets written into
the next generation's starting priors. In human cognition and cultural evolution,
this is precisely the question that \textit{is} answered: credibility-enhancing
displays, prestige-biased transmission, and community-level attention allocation
are the mechanisms by which prior preferences over outcomes are shaped and
reproduced~\citep{henrich2009CREDs}. Any formal account of norm learning that
cannot speak to these mechanisms is incomplete at the level of mechanism, even
if it is complete at the level of equilibrium description.

The stakes of this incompleteness become acute when AI systems participate in
human normative life. \citet{kulveit2025GradualDisempowerment} argue that
incremental AI integration can gradually disempower human populations not through
any adversarial takeover but through structural drift: as AI systems become
load-bearing participants in economic, cultural, and institutional life, the
implicit alignment between human interests and system behaviour that comes from
human participation being generative is progressively eroded. Under the equilibrium
framing of norms, AI systems that sit at the behavioural equilibrium look
unproblematic. But if norms supervene on shared cognitive architecture --- if
their stability is grounded in the precision-shaping dynamics through which a
community installs prior preferences in its members --- then AI systems that
participate only at the policy level, without contributing to the substrate that
makes norms credible, can hollow out the normative structure even while the
behaviour looks intact. The installation question is therefore not merely a
technical gap in formal models; it is the micro-level mechanism through which
gradual disempowerment operates.

\begin{figure}[t]
\centering
\begin{tikzpicture}[
  node distance = 2.4cm and 3.2cm,
  box/.style = {
    draw, rectangle, rounded corners=3pt,
    minimum width=2.4cm, minimum height=0.9cm,
    font=\small, align=center
  },
  missing/.style = {
    draw=gray!60, rectangle, rounded corners=3pt,
    minimum width=2.4cm, minimum height=0.9cm,
    font=\small, align=center, dashed, fill=gray!10
  },
  arr/.style  = {-Stealth, thick},
  marr/.style = {-Stealth, thick, dashed, gray!70}
]
  %% Main horizontal chain
  \node[box] (world)   {World};
  \node[box, right=of world]  (policy)  {Policy $\policy$};
  \node[box, right=of policy] (action)  {Action};

  %% Off-axis community node
  \node[missing, above=1.5cm of policy] (community) {Community};

  %% Horizontal arrows
  \draw[arr]  (world)  -- (policy);
  \draw[arr]  (policy) -- (action);

  %% Missing arrow: community -> prior layer (annotated)
  \draw[marr] (community) -- node[right, font=\footnotesize, text=gray!80]
              {missing / proposed} (policy);

  %% Label under the policy node to indicate the prior layer
  \node[font=\footnotesize, below=0.15cm of policy, text=gray!60]
        {prior $\C$};

\end{tikzpicture}
\caption{The missing arrow: existing cooperative-AI norm models operate on the
horizontal chain (world $\to$ policy $\to$ action) and leave the community-to-prior
channel exogenous (dashed). This paper's contribution is to formalise the shaded
arrow: how community-level precision dynamics install an agent's prior preferences
over outcomes ($\C$), constituting the cognitive substrate of norms.}
\label{fig:missing-arrow}
\end{figure}

\subsection*{Differentiation from Albarracin et al.\ (2025)}

\citet{albarracin2025SocialPower} formalise social power as the differential
capacity to set the policy-precision parameter ($\prc$) in others' generative
models. In their account, a socially powerful agent is one who can raise or
lower how confidently another agent pursues a given policy --- shaping the
effective temperature of the policy distribution without necessarily altering
what the agent considers valuable in the first place. This is a significant
contribution: it gives power a formal active-inference definition, it engages
Bourdieu's concepts of habitus, field, and capital, and it shows how
Foucauldian disciplinary mechanisms can be read as precision-modulation
operations. Our move is one layer up, on the \textit{content} of prior
preferences ($\C$) --- that is, what an agent comes to consider valuable
in the first place, rather than how confidently it pursues a given policy.
Both layers are causally involved in the cultural reproduction of norms:
a community that controls the policy-precision parameter $\prc$ in its
members (the Albarracin layer) can enforce compliance, but a community that
controls the content of $\C$ (our layer) shapes what counts as desirable
before any enforcement becomes necessary. The prior-preference content layer
is, we argue, the substrate that cooperative-AI norm theory currently lacks.
Albarracin et al.\ provide the complementary half of the story; this paper
provides the other half, and together they constitute a two-layer account of
normativity in active-inference agents. The dynamics by which epistemic
communities stabilise or fragment around shared $\C$-vectors are developed
in~\citet{albarracin2022EpistemicCommunities}, which we build on directly
in Section~\ref{sec:formal}.

\subsection*{Roadmap}

The remainder of the paper develops the missing arrow in three steps.
Section~\ref{sec:cultural-evolution} establishes the cultural-evolution backdrop,
reviewing the mechanisms --- prestige-biased transmission, credibility-enhancing
displays~\citep{henrich2009CREDs}, and community-level attention allocation ---
through which prior preferences are differentially shaped and reproduced across
agents.
Section~\ref{sec:formal} develops the formal account: we define norms as
stabilised regions of high-precision shared $\C$, characterise the installation
channel as a community-level precision-shaping process, and derive stability
conditions for the resulting attractor landscape.
Section~\ref{sec:simulation} proposes a minimal simulation instantiating the
account in a pymdp multi-agent environment, designed to exhibit the bifurcation
between norm-convergence and echo-chamber fragmentation as a function of
the precision dynamics.


\section{Where Values Come From: Prestige, CREDs, and the Cultural-Evolution
Substrate}\label{sec:cultural-evolution}

Humans are, in a fundamental sense, obligate cultural learners: the overwhelming
share of what any individual knows, wants, and treats as worth attending to was
acquired from other people rather than rediscovered from scratch
\citep{henrichGilWhite2001Prestige}.
This dependence runs deeper than the transmission of practical skills or factual
beliefs.
What counts as a \emph{value} — which outcomes are desirable, which actions are
admirable, which objects are sacred or disgusting — is itself downstream of what
gets attended to, imitated, and rehearsed in a community's daily life.
A child who grows up watching elders fast, or observing neighbours donate
conspicuously to shared projects, or noticing which adults attract the gaze and
proximity of many others, acquires not just a collection of beliefs but a
structured landscape of salience: some things matter a great deal, others barely
register.
The cultural-evolution tradition, crystallised in \citet{henrichGilWhite2001Prestige},
has spent the past quarter-century documenting the population-level mechanisms
that produce and maintain these landscapes.
This section surveys that account in enough detail to reveal where a formal
bridge to predictive processing is both possible and necessary.

\subsection*{Prestige as a coordination device for collective attention}

The central conceptual move in \citet{henrichGilWhite2001Prestige} is the
distinction between \emph{prestige} and \emph{dominance} as two structurally
different routes to social influence.
Dominance hierarchies, which appear across social primates, are sustained by
intimidation: individuals defer to dominant animals to avoid costs.
Prestige, by contrast, is \emph{freely conferred deference}.
Copiers voluntarily attend to, sit near, gift, and imitate high-prestige models
because doing so raises their own fitness by providing access to superior
information or skill.
The evolutionary logic is straightforward: when acquiring complex cumulative
culture requires accurately identifying good models, individuals who can cheaply
infer model quality by reading others' deference choices will outcompete those
who must independently assess each candidate.
The community's existing attention distribution therefore functions as a
\emph{second-order prior} over whom it is worth attending to — a public signal
that learners consult before committing their own attentional resources.

\citet{henrichChudekBoyd2015BigMan} formalise this logic in a
culture--gene coevolutionary model of ``big man'' leadership.
In small-scale egalitarian societies, the big man acquires followers not through
coercion but through accumulated demonstrations of generosity, oratory, and
successful co-ordination.
Critically, followers copy big-man behaviours (including prosocial ones), which
in turn creates selection pressure for big men who are genuinely prosocial —
because their only currency is the continued attention and copying their conduct
commands.
The model shows that this prestige channel produces phenotype correlation between
leaders and followers and generates prosocial norms as stable downstream
outcomes: the attention asymmetry is not merely a readout of pre-existing value
distributions but actively reshapes them.
Egalitarian levelling mechanisms — gossip, mockery, ostracism, and in extremis
execution — function not as alternatives to prestige dynamics but as its
complement: they \emph{withdraw} community attention from would-be dominants,
enforcing the condition under which prestige (rather than coercion) remains the
operative currency.
Flat societies are densely attention-managed, not attention-neutral.

The clearest experimental demonstration of how this second-order attention
mechanism works at the level of individual learning is \citet{chudek2012BystanderGaze}.
In a series of studies with 3--4-year-old children, a purely \emph{attentional}
cue — ten seconds of differential bystander gaze toward one of two adult
demonstrators, with no other information about their competence or history — more
than doubled the probability that children subsequently copied that demonstrator
on a novel task.
Crucially, the effect \emph{generalised across domains}: prestige established
through others' gaze in one context (e.g., artifact use) transferred to
unrelated copying choices (e.g., food preferences).
This domain-generality is theoretically decisive.
A content-specific prior — ``this person knows about artifacts'' — would not
transfer to food choices; the prestige cue would stay in the domain where it was
established.
What does transfer is a \emph{precision signal}: ``this person's behaviour is
worth attending to, period.''
In predictive-processing terms (though the authors do not use this vocabulary),
the bystander gaze cue behaves like a generalised precision parameter rather than
a domain-restricted content prior.
It raises the weight accorded to the high-prestige model's actions across the
entire space of observable behaviours, without specifying what those actions will
contain.
Section~\ref{sec:formal} formalises this observation precisely.

\subsection*{CREDs and the costliness logic}

Once language evolved to allow the cheap transmission of beliefs, cultural
learners faced a new inference problem: which professed values does a model
\emph{actually hold}, and which are mere verbal performance?
\citet{henrich2009CREDs} argues that natural selection favoured learners who
weighted models' expressed beliefs by the costliness of their observable
commitment to those beliefs — what he calls Credibility-Enhancing Displays
(CREDs).
Fasting, painful initiations, celibacy, martyrdom, and public sacrifice are
actions whose costs make sense only if the agent genuinely holds the associated
belief or value; they are costly signals that cannot be cheaply faked.
Learners who copied the values of agents who performed CREDs systematically
acquired the values of genuinely committed models, rather than those of
convenient verbal performers.

The fit with precision-weighted inference is near-direct, and Section~\ref{sec:formal}
develops it explicitly.
Informally: without CREDs, a model's verbal claim about a value provides
low-precision evidence — it could reflect genuine commitment or could be
strategic posturing.
A CRED raises the precision of the channel from observed-other-behaviour to
one's own value-prior.
High-CRED demonstrations are high-confidence messages about what the transmitter
actually values; low-CRED environments produce noisy, uncertain updates.
Empirical work by Lanman and Buhrmester \citep{lanmanBuhrmester2017CREDs} has
confirmed this logic: childhood exposure to parental CREDs predicts adult theism
more robustly than competing accounts based on cognitive style or existential
anxiety, suggesting that what gets installed is not a bare belief but a
\emph{calibrated prior} whose strength tracks the evidential quality of the
exposure.
The CRED mechanism and the prestige mechanism are thus complementary channels in
the same installation system: prestige cues direct \emph{which} models receive
attention; CREDs modulate \emph{how much confidence} to accord to those models'
expressed commitments.

\subsection*{Attentional infrastructure is itself culturally inherited}

A natural objection at this point is that while the \emph{contents} of attention
may vary culturally, the \emph{mechanisms} that structure attention — what kinds
of agents and events are even candidates for imitation — must be
evolutionarily fixed.
\citet{heyes2018CognitiveGadgets} mounts a sustained argument that this
objection is wrong.
Her ``cognitive gadgets'' thesis distinguishes a domain-general genetic
``starter kit'' of social attentional biases (sensitivity to faces, voices, and
social rewards) from the higher-order mechanisms of selective social learning
— including imitation, mind-reading, and norm compliance — which are themselves
culturally assembled during development.
The decisive piece of evidence for present purposes: there is empirical support
for the claim that humans can \emph{socially learn attentional biases},
specifically by learning to attend more to some agents than others through the
mechanism of following third-party gaze.
That is, the same bystander-gaze cue that Chudek et al.\ use as an experimental
manipulation of prestige is also the mechanism by which children acquire durable
attentional dispositions — what to find exemplary, whom to watch, which features
of the social environment carry signal.
The salience landscape is not merely populated through cultural transmission; its
\emph{topography} — which gradients channel attention up rather than down — is
itself inherited.

\citet{heyes2024NormPsychology} extends this framework to norm psychology
specifically.
Rather than treating norm-following as a genetically fixed motivational module
(``normative competence as instinct''), she argues that normative competence is a
culturally assembled cognitive gadget built atop domain-general attentional and
motivational capacities.
Crucially, what counts as a norm — which regularities in others' behaviour are
picked up as prescriptive rather than merely descriptive — is not given in
advance but is itself acquired through the same prestige-and-gaze mechanisms
that structure learning more generally.
The implications for any account of value transmission are significant: there is
no innate, domain-general ``norm receptor'' that picks out which social patterns
to treat as obligatory.
Instead, the very categories through which norms are perceived and reproduced
are part of what gets transmitted.
This recursive structure — attentional dispositions shaping what gets copied,
which in turn shapes attentional dispositions in the next generation — is
precisely the loop this paper seeks to formalise.

\subsection*{The ceiling of the cultural-evolution account}

The literature surveyed in this section provides a rich \emph{empirical} story
about value transmission: which individuals attract community attention, under
what conditions, through what cues, with what consequences for the behavioural
repertoires of those who copy them.
The account is anchored in experimental data \citep{chudek2012BystanderGaze},
population-genetic formalisms \citep{henrichChudekBoyd2015BigMan,egoziRam2024PrestigeFormal},
and developmental evidence \citep{heyes2018CognitiveGadgets,heyes2024NormPsychology}.
What this literature does \emph{not} provide is a vocabulary that connects these
dynamics to the computational architecture of the individual learner: why does
bystander gaze function as a domain-general precision signal rather than a
domain-specific content cue?
What is happening, formally, when a CRED raises the confidence with which a
value-prior is adopted?
How do community-level patterns of attention crystallise into the stable
configurations we call norms?
These questions ask for a level of mechanistic description that cultural
evolution's largely informal vocabulary does not supply.
The next section takes up precisely this gap, mapping prestige- and
CRED-mediated attentional dynamics onto the formal language of
precision-weighted predictive processing, and showing how the loop that closes
from community priors through attention and copying back to community priors can
be given a rigorous active-inference rendering.


\section{Values as Salience Priors: A Formal Account}\label{sec:formal}

%% ~3500 words / ~4–5 pages. Heaviest, most technical section.
%% Anchors: friston2015ActiveInferenceEpistemic, hesp2021DeeplyFelt,
%%   feldmanFriston2010Attention, parrFriston2017Attention,
%%   parrFriston2019AdvancedAtt, parviziWayne2024PreferencesAttention,
%%   limanowskiFriston2018Dark, torresan2025ShapeC, millidge2021EFE,
%%   ransom2020Cognition, veissiere2020TTOM, albarracin2022EpistemicCommunities,
%%   albarracin2025SocialPower, constant2019Regimes, heins2024CollectiveSurprise,
%%   heins2025AsOneAndMany, vasil2020SameKind, henrichGilWhite2001Prestige,
%%   henrich2009CREDs, chudek2012BystanderGaze, koster2022SillyRules.
%% Defensive: millidge2021EFE in a footnote.

\subsection{Attention is precision; salience is expected information gain}
\label{sec:attention-vs-salience}

The first step is to refuse a conflation that even active-inference researchers
routinely commit. \citet{feldmanFriston2010Attention} establish the load-bearing
identification: \emph{attention} is the inference of precision $\Prec$ on
prediction errors --- a gain on the likelihood mapping that determines which
channels of current sensory evidence dominate posterior updating. Higher
precision at a level of the hierarchy means a louder error signal at that
level, and louder errors do more work in revising beliefs. Salience is a
different quantity. \citet{parrFriston2017Attention,parrFriston2019AdvancedAtt}
draw the distinction with full formality: salience is the \emph{expected
information gain} associated with a candidate sample of the world ---
prospective and policy-conditional, not retrospective and stimulus-bound. To a
first approximation, attention is precision over the data the agent is
\emph{currently} consuming; salience is a forward-looking quantity that scores,
for each candidate next observation, how much the agent's posterior would
contract were that observation actually obtained.\footnote{We will write the
salience operator $\salience$ informally as ``expected information gain'' and
will not commit to the full expected-free-energy decomposition. The logical
content of our argument requires only that preferences and information gain
appear additively in a single objective; we do not require that the
expected-free-energy functional be uniquely derivable in the sense of
\citet{millidge2021EFE}, whose foundational concerns about EFE we acknowledge
but do not need to resolve here.} Salience is therefore not a property of the
distal world. It is a property of the prior-shaped uncertainty landscape an
agent carries around with it.

The framing matters because, once one allows that deep priors set the precision
weights distributed across the hierarchy, salience inherits its shape from
those priors. \citet{parviziWayne2024PreferencesAttention} makes the point most
forcefully: under active inference, the endogenous--exogenous attention
dichotomy collapses. There is no precision-free signal entering the cognitive
system. What looks like bottom-up, stimulus-driven capture is in fact the
allocation of precision against prediction errors generated by deep generative
expectations the agent already holds. The corollary, due to
\citet{limanowskiFriston2018Dark}, is that phenomenal transparency and opacity
are themselves precision phenomena --- when precision on a level is high and
stable, the agent ``sees through'' it; when precision is low or unsettled, the
level itself becomes available to introspection. Together, these results
license the move at the heart of this paper: \emph{deep priors --- including
the agent's preferences --- shape what becomes salient}. \citet{ransom2020Cognition}
press the strongest empirical objection to attention-as-precision: affect-biased
attention orients the system toward emotionally significant stimuli even when
their feature-level precision is low. We engage this directly because it is
not in fact a counter-example. Their data are consistent with a deep-priors-set-
precision rebuttal: emotionally significant stimuli are exactly the ones for
which the agent has high-precision priors over their value-relevance, and
those high-precision deep priors propagate down the hierarchy as gain
modulation on the lower-level error signals that constitute affect-biased
attention. The Ransom et al.\ result, in other words, is evidence \emph{for}
the framing this paper develops, provided one is willing to grant that values
sit upstream of sensory precision-setting --- which is exactly the move we now
make.

\subsection{Values as prior preferences $\C$}
\label{sec:values-as-C}

Active inference furnishes a single object that does the work of
``preferences,'' ``goals,'' ``values,'' and ``utility'' in a unified register.
The prior preference distribution $\C$ encodes the agent's expectations about
\emph{which observations} it will tend to obtain. Crucially, these are
expectations in the ordinary Bayesian sense: $\C$ is a prior over outcomes,
and the agent acts in such a way as to make its sensory experience consistent
with this prior \citep{friston2015ActiveInferenceEpistemic}. ``Wanting'' an
outcome and ``expecting'' it become two faces of one variational quantity.
\citet{friston2015ActiveInferenceEpistemic} establish the formal apparatus that
puts pragmatic value (the expected log-probability of preferred outcomes) and
epistemic value (expected information gain about hidden states) into a single
additive objective; preferences and salience appear as commensurable terms,
not as separate cognitive systems. Importantly, our argument requires only the
\emph{decomposition} --- preferences and information are summed in a single
currency --- and not the strong claim that this currency is uniquely fixed by
expected free energy.

The bridge from values-as-priors to precision is supplied by
\citet{hesp2021DeeplyFelt}. They formally identify valence with the
\emph{expected precision of the action model}: a positively valenced state of
affairs is one in which the agent expects its model to reliably support
goal-directed action; a negatively valenced one is one in which it expects
that model to be unreliable. ``Caring about'' an outcome, on this account, is
\emph{having a high-precision prior expectation of obtaining it}. This is
the move that turns the abstract notion of $\C$ into a precision-bearing
object. A community that shares values is a community whose members hold
similarly shaped, similarly precise $\C$'s --- and, by the result of
\S\ref{sec:attention-vs-salience}, a community whose members find the same
features of the world salient.

\begin{definition}[Values as prior preferences]
\label{def:values}
An agent's \emph{values} over outcomes are encoded as the prior preference
distribution $\C$ in its generative model. The valence of an outcome $\obs$ is
$\log p(\obs \mid \C)$; high-precision $\C$ expresses sharp values and
low-precision $\C$ expresses diffuse values. Equivalently, a value is a
high-precision prior expectation of observing a class of outcomes, and the
felt strength of a value is the precision of that expectation.
\end{definition}

This definition has two immediate consequences. First, it dissolves the
intuitive contrast between ``what an agent wants'' and ``what an agent
expects'': both are encoded by $\C$, and the dual reading is constitutive of
the active-inference account, not a translation artefact. Second, it makes
``valuing'' a graded property: there is no all-or-nothing distinction between
having and not having a value, only the continuous question of how peaked
$\C$ is over an outcome class. \citet{constant2019Regimes} extend this picture
to the social register by introducing \emph{deontic value}: a third value
type alongside epistemic and pragmatic, expressed as policy priors cued by
environmentally and socially supplied scripts. Deontic cues (a red traffic
light, a raised eyebrow, an averted gaze) emerge from repeated collective
behaviour and reify cultural-group preferences as Bayes-optimal action
defaults. Deontic priors are policy-level priors --- priors over what to do
given a cue --- whereas $\C$ is an outcome-level prior. The two layers
interact: outcome-level preferences shape which policies look pragmatically
valuable, and deontic policy priors shape which policies get selected without
explicit deliberation. Section~\ref{sec:installation} formalises how the
outcome-level layer, $\C$, is itself acquired from a community.

\subsection{The $\C$-installation channel}
\label{sec:installation}

The genuinely new piece of formal machinery this paper supplies concerns how
$\C$ is \emph{installed}. Existing accounts treat $\C$ as exogenously
specified by the modeller. We propose instead that $\C$ is acquired from a
community, by a precision-modulated channel from observed-other-behaviour into
the agent's own preference distribution. Schematically:
\begin{equation}
\label{eq:channel}
\text{community behaviour} \;\xrightarrow{\;\;\;\salience \cdot \prc\;\;\;}\;
\text{updated } \C,
\end{equation}
where $\prc$ is a scalar channel precision and $\salience$ is the salience
allocated to the observed behaviour. The channel is bidirectional in the
deeper sense that its input $\salience$ is itself determined by the agent's
current $\C$ (via \S\ref{sec:attention-vs-salience}); the agent's existing
preferences shape which conspecific behaviours it samples and weighs, and the
samples it weighs shape its updated preferences. Equation~\eqref{eq:channel}
is therefore not a one-shot transmission rule but the elementary step of a
fixed-point dynamic which we describe in \S\ref{sec:loop}.

The empirically novel claim is that prestige cues and credibility-enhancing
displays --- the documented mechanisms of human cultural transmission ---
function as modulators of the channel precision $\prc$, not as content
priors. \citet{henrichGilWhite2001Prestige} establish prestige as
freely-conferred deference: high-prestige individuals are not coerced into
attention-monopoly but selected for it by observers as a heuristic for
locating skilled or successful models. \citet{henrich2009CREDs} argues that
costly displays of commitment to beliefs and norms function as honest signals
that increase observers' willingness to acquire those beliefs and norms.
Both mechanisms have been treated in cultural-evolutionary theory as
\emph{biases on what to copy}; we propose they are better treated as
modulators of \emph{how strongly to update} on observation, that is, as gains
on $\prc$. The empirical anchor for this reframing is
\citet{chudek2012BystanderGaze}, whose experimental finding has been
under-exploited in the active-inference literature. Children observed
multiple potential models, with bystanders' gaze biased toward one of them.
Children subsequently preferred to copy the bystander-attended model across
\emph{multiple} content domains --- not just the domain in which the
bystander gaze occurred. This generality is the formal signature of a
precision parameter rather than a content prior: a content prior would
predict transfer only to similar contents, while a generalised precision
parameter predicts transfer across the whole channel through which content
flows. Bystander gaze, on this reading, is a CRED-like cue whose mechanistic
effect is to elevate $\prc$ on the channel from the attended model into the
observer's $\C$.

This reading must be distinguished from two adjacent accounts.
\citet{veissiere2020TTOM}'s Thinking Through Other Minds framework treats
culture as installing \emph{regimes of attention} --- shared styles of
allocating attentional resources that mark some channels as reliable. TTOM
operates at the level of shared generative-model structure and the precision
weights that style channel reliability; it is silent on the content of $\C$.
\citet{vasil2020SameKind} argue that human communication is enabled by a
deep prior that one's interlocutor is the ``same kind of creature'' as
oneself, licensing rapid alignment of generative models. This is also a
structural prior, not a content prior. We do not take issue with either
account; we extend them. The TTOM regime determines \emph{whose} $\C$ a
learner samples; the same-kind prior makes that sampling possible at all;
prestige and CRED cues then determine \emph{how much} the sampled
behaviour is allowed to modify the learner's $\C$. The channel of
Equation~\eqref{eq:channel} formalises the third move, which the prior
literature gestures at without naming. Differentiation from
\citet{albarracin2025SocialPower} requires special care, since their account
is the closest existing neighbour. They formalise social power as the
differential capacity to set policy-precision (often denoted $\alpha$) in
others' generative models. Power, on their account, sits at the
\emph{policy-precision} layer: empowered agents anchor others' policy priors
to their own outputs and benefit from the resulting attention asymmetry. Our
account sits at the \emph{prior-preference content} layer:
$\C$-installation modifies what others come to want, not which policies they
deploy in service of pre-given wants. The two layers are coupled --- a
community whose $\C$ has been installed by prestige cues will in addition
exhibit asymmetric policy-precision around the prestige-bearer --- but they
are formally separable, and our claim is that the $\C$-content channel is
the one cultural evolution has been describing for half a century without
naming it as such.

\subsection{The fixed-point loop}
\label{sec:loop}

The four moves --- values are $\C$ (\S\ref{sec:values-as-C}); $\C$ shapes
salience via precision (\S\ref{sec:attention-vs-salience}); salience directs
copying through the prestige/CRED-modulated channel (\S\ref{sec:installation});
copying installs $\C$ in the next-generation agent --- close into a loop.
Figure~\ref{fig:loop} renders the loop graphically.

\begin{figure}[t]
\centering
%% Manual TikZ positioning is used here for pdfLaTeX compatibility.
%% A `graphs`+`graphdrawing` (layered/circular) version is available but
%% requires LuaLaTeX; the manual layout below produces a clean circular
%% diagram under either compiler.
\begin{tikzpicture}[
  >=Stealth,
  node distance=18mm,
  every node/.style={font=\small},
  prior/.style={draw=blue!70!black, fill=blue!8, thick, rounded corners,
                inner sep=4pt, align=center, minimum width=24mm,
                minimum height=10mm},
  precision/.style={draw=orange!80!black, fill=orange!10, thick, rounded corners,
                    inner sep=4pt, align=center, minimum width=24mm,
                    minimum height=10mm},
  social/.style={draw=red!70!black, fill=red!8, thick, rounded corners,
                 inner sep=4pt, align=center, minimum width=24mm,
                 minimum height=10mm},
  arr/.style={->, thick}
]
  %% Five nodes arranged on a circle
  \node[prior]     (commC)   at ( 90:36mm) {community $\C$};
  \node[precision] (sal)     at ( 18:36mm) {salience\\landscape};
  \node[social]    (att)     at (-54:36mm) {attention \&\\copying};
  \node[social]    (channel) at (234:36mm) {prestige / CRED\\channel ($\prc$)};
  \node[prior]     (instC)   at (162:36mm) {installed $\C$\\(next agent)};

  %% Cycle arrows
  \draw[arr, blue!70!black]   (commC)   -- (sal);
  \draw[arr, orange!80!black] (sal)     -- (att);
  \draw[arr, red!70!black]    (att)     -- (channel);
  \draw[arr, red!70!black]    (channel) -- (instC);
  \draw[arr, blue!70!black]   (instC)   -- (commC);

  %% Inner gloss
  \node at (0,0) {\itshape fixed-point};
\end{tikzpicture}
\caption{The values-as-salience-priors loop. A community's prior preferences
($\C$) determine the salience landscape, which directs attention and copying
via prestige- and CRED-modulated channels; the resulting installed $\C$ in
next-generation agents reproduces (or shifts) the community-level prior.
Colours: blue = prior content; orange = precision (salience); red = social
channel.}
\label{fig:loop}
\end{figure}

The loop has three qualitatively distinct dynamical regimes. \emph{Convergence}
obtains when channel precision $\prc$ is high enough and prestige consensus
is sufficient that successive generations install $\C$'s closely matching
the community-level distribution; the loop reproduces the prior with high
fidelity, and norms are robust. \emph{Bifurcation} occurs when the loop
factorises across sub-populations whose internal $\prc$ exceeds the
cross-population $\prc$: \citet{albarracin2022EpistemicCommunities} demonstrate
this regime explicitly with their epistemic-community simulations, in which
shared priors over source reliability lock sub-communities into echo-chamber
attractors that resist cross-community correction. \emph{Dissolution} occurs
when $\prc$ collapses across the whole community --- when prestige cues lose
their force, CRED bearers are absent or distrusted, or the salience
landscape becomes too noisy to support coherent copying --- and the
community-level $\C$ flattens, as norms lose precision and content drifts.
A subtle but important point follows from \citet{heins2025AsOneAndMany}: the
community-level $\C$ is \emph{not} simply the average of individuals' $\C$'s.
Group-level generative models of a coupled collective have non-trivial
emergent structure, and the equilibrium behaviour of the loop in
Figure~\ref{fig:loop} is a property of the coupled system, not a moment of
the individual-level distribution. We adopt this caution: when we speak of
``the community's $\C$,'' we mean the aggregate that emerges from the loop
dynamics, not a population mean.

\subsection{Norms as stabilised regions of high-precision shared $\C$}
\label{sec:norms}

The loop dynamics give us a precise formulation of what norms are. A norm is
not a rule, an equilibrium, or a piece of representational content per se;
it is a region of outcome space where the community's $\C$ has both high
mass and high precision, stabilised by the loop of
\S\ref{sec:loop}.

\begin{definition}[Norm]
\label{def:norm}
A \emph{norm} in a community is a region $\norm \subseteq \mathcal{O}$ of
outcome space such that the community-level prior preference distribution
has both high mass on $\norm$ and high precision around $\norm$, where the
precision is sufficient that observed deviations from $\norm$ generate
prediction errors strong enough to drive both perception (the violation
\emph{feels} wrong) and action (correction, sanction, or attention-withdrawal).
\end{definition}

The definition recovers and explains three otherwise puzzling features of
norms. First, it explains why \emph{violations of norms feel surprising in a
way that violations of mere regularities do not}: high-precision priors over
$\C$ generate large prediction errors when violated, and high-precision
errors are exactly what valenced surprise consists of
\citep{hesp2021DeeplyFelt}. Second, it explains the ``silly rules'' result of
\citet{koster2022SillyRules}, where adding arbitrary content-free taboos
into a multi-agent reinforcement-learning environment improves agents'
acquisition of welfare-relevant norms. Under our account, what silly rules
do is pump up the precision of the meta-prior ``a normative regime is in
operation'' without committing to particular content; once that precision is
elevated, agents acquire welfare-relevant norms faster because the
$\C$-installation channel is opened wider. The mechanism Köster et al.\
identify is exactly the precision-modulator move of \S\ref{sec:installation},
written in MARL form. Third, the definition makes contact with the
deontic-value apparatus of \citet{constant2019Regimes}: norms-as-regions of
high-precision $\C$ generate the environmental and social cues that, in a
mature normative culture, become deontic priors over policies.
Outcome-level $\C$ and policy-level deontic priors thus stand in a
relationship that the active-inference literature already supports: deontic
cues short-circuit pragmatic deliberation when the loop has stabilised the
underlying $\C$, allowing norm-following to occur as Bayes-optimal action
under environmentally cued priors. \citet{heins2024CollectiveSurprise}
furnish the closest existing demonstration that norm-like collective
patterns can emerge from agents minimising variational free energy under
coupled generative models, without explicit social-force or utility rules.
Our claim is that the same logic, applied to the prior-preference content
layer rather than the kinematic state layer, recovers the cultural-norms
regime.

\subsection{The shape of $\C$: sharp, broad, and shaped}
\label{sec:shape-of-C}

Norms differ in more than location. They differ in the \emph{shape} of $\C$
they instantiate. \citet{torresan2025ShapeC} provide the empirical
demonstration that drives this subsection: in a grid-world navigation task,
they compare four parameterisations of $\C$ --- hard goals (delta-like
peaked priors), soft goals (broad priors), and either of these with or
without intermediate goal shaping. Their finding has direct consequences for
our account of norms.

\begin{figure}[t]
\centering
%% Three-panel PGFPlots figure laid out via \pgfplotsset groupplots-style
%% positioning would require an extra \usepgfplotslibrary{groupplots} in the
%% preamble; we instead use three independent axis environments arranged
%% horizontally inside a single tikzpicture, which works under the existing
%% pgfplots load in main.tex without further dependencies.
\begin{tikzpicture}
  \pgfplotsset{
    every axis/.append style={
      width=4.6cm, height=3.8cm,
      xmin=-4, xmax=4, ymin=0, ymax=0.85,
      xtick={-3,0,3}, ytick=\empty,
      axis lines=left,
      every axis label/.style={font=\footnotesize},
      every axis title/.style={font=\small},
    }
  }
  \begin{axis}[
    name=panA, at={(0,0)}, anchor=south west,
    xlabel={\small outcome $\obs$}, ylabel={\small $p(\obs\mid\C)$},
    title={\small (A) sharp $\C$}
  ]
    \addplot[thick, blue!70!black, smooth, samples=200, domain=-4:4]
      {0.8*exp(-((x-0)^2)/(2*0.25^2))};
  \end{axis}
  \begin{axis}[
    name=panB, at={(panA.south east)}, anchor=south west, xshift=8mm,
    xlabel={\small outcome $\obs$},
    title={\small (B) broad $\C$}
  ]
    \addplot[thick, blue!70!black, smooth, samples=200, domain=-4:4]
      {0.28*exp(-((x-0)^2)/(2*1.4^2))};
  \end{axis}
  \begin{axis}[
    name=panC, at={(panB.south east)}, anchor=south west, xshift=8mm,
    xlabel={\small outcome $\obs$},
    title={\small (C) shaped $\C$}
  ]
    \addplot[thick, blue!70!black, smooth, samples=200, domain=-4:4]
      {0.6*exp(-((x-0)^2)/(2*0.4^2))
       + 0.2*exp(-((x-2)^2)/(2*1.0^2))
       + 0.18*exp(-((x+2.2)^2)/(2*1.1^2))};
  \end{axis}
\end{tikzpicture}
\caption{The shape of $\C$ governs the exploration--exploitation tradeoff
and the robustness of installed values. (A) Sharp $\C$: hard goals; fast
convergence; fragile to perturbation. (B) Broad $\C$: soft goals; preserved
exploration; slow drift. (C) Shaped $\C$: a peak with structured shoulders
on either side --- most performant per \citet{torresan2025ShapeC}; a hybrid
of pragmatic and epistemic value.}
\label{fig:shape}
\end{figure}

A sharp $\C$ (Figure~\ref{fig:shape}A) is a delta-like preference over a
narrow outcome region. Sharp priors push the expected-free-energy objective
toward pragmatic value at the expense of epistemic value: agents converge
fast on the preferred outcome but under-sample alternatives, fragility
follows, and small environmental perturbations that move the optimal action
off the prior's peak are not detected because the precision penalty on
sampling alternatives is large. A broad $\C$ (Figure~\ref{fig:shape}B) is a
diffuse preference covering many outcomes. Broad priors preserve epistemic
value but achieve slower convergence on any particular target; in the limit,
a perfectly flat $\C$ collapses the pragmatic term entirely, and the agent
becomes a pure information-seeker. A shaped $\C$ (Figure~\ref{fig:shape}C)
combines a sharp peak on a target outcome with structured residual mass on
neighbouring outcomes --- the goal-shaping regime
\citet{torresan2025ShapeC} identify as most performant. Shaped $\C$
preserves the pragmatic gradient toward the preferred outcome while
maintaining enough probability on alternatives to support continued
exploration.

The translation to norms is direct. A community whose shared $\C$ is sharp
is a community of strongly aligned values that converges fast and resists
perturbation poorly --- a fragility property of normative coherence that is
otherwise hard to articulate. A community whose shared $\C$ is broad
preserves exploration of normative alternatives but coordinates poorly. A
community whose shared $\C$ is shaped --- a peak of strong commitments
flanked by structured latitude --- combines the convergence properties of
the first with the robustness of the second. This maps onto a familiar
cultural-evolutionary observation: healthy traditions tend to combine sharp
core commitments with broad peripheral latitude, and the precision of $\C$
is graded across content rather than uniformly high. Definition~\ref{def:norm}
must be read in this light: a norm is a high-precision \emph{region}, not a
high-precision \emph{point}, and the shape of the surrounding gradient is a
substantive property of the norm with measurable consequences for the
loop's robustness.

\subsection{Synthesis}
\label{sec:formal-synthesis}

The four moves of this section close the gap identified in \S\ref{sec:cooperative-ai}.
Values are formally identified with the prior preference distribution $\C$
(Definition~\ref{def:values}); $\C$ is installed by a
prestige- and CRED-modulated precision channel from observed-other-behaviour
into the agent's own preferences (\S\ref{sec:installation}); norms are
stabilised regions of high-precision shared $\C$ in a community
(Definition~\ref{def:norm}), held in place by the fixed-point loop of
Figure~\ref{fig:loop}; and the shape of $\C$ governs the exploration--exploitation
balance and therefore the robustness of installed values
(\S\ref{sec:shape-of-C}). The community-to-$\C$ arrow that the literature
described without naming is now formal. It is also the arrow on which the
cooperative-AI implications of Section~\ref{sec:simulation} turn: a system that
mediates the prestige- and CRED-bearing channel mediates the precision of
the cultural-transmission machinery itself, with consequences for the
stability of human normative orders that we take up next.


\section{A Proposed Simulation: Norm Formation under Prestige-Biased Precision}%
\label{sec:simulation}

The formal account in Section~\ref{sec:formal} yields a set of testable predictions about
regime transitions in a community of active-inference agents: as the precision
gain on inter-agent observation channels rises, the community should traverse
states of heterogeneous individual priors, broad consensus, bifurcated
echo-chamber attractors, and finally degenerate fixation on a single agent's
prior. We sketch here a minimum viable simulation that operationalises these
predictions. The base platform is the pymdp epistemic-communities setup of
\citet{albarracin2022EpistemicCommunities}, which already demonstrates that
self-confirming sampling under shared priors produces stable echo chambers
without any malign design. We extend that setup with an explicit
\emph{prestige-cue mechanism} that modulates the precision on inter-agent
observation channels, allowing us to sweep through the predicted regime
sequence in a single parameter.

\subsection*{Setup}

We consider $N \approx 20$ agents, each running a discrete-state active-inference
loop implemented in pymdp. Each agent $i$ maintains a generative model with a
shared discrete outcome space $\mathcal{O}$ and a privately initialised prior
preference vector $\C_i \in \mathbb{R}^{|\mathcal{O}|}$ (the log-preference, in
the sense of \citealt{torresan2025ShapeC}).  Agents observe each other's most
recent actions through a precision-weighted observation channel: the
\emph{effective channel precision} from agent $j$ to agent $i$ is
\begin{equation}
  \prc_{ij} \;=\; \prc_0 \cdot f(\rho_{ij}),
  \label{eq:prestige-precision}
\end{equation}
where $\prc_0 > 0$ is a baseline precision, $\rho_{ij} \in [0,1]$ is a
prestige cue, and $f : [0,1] \to \mathbb{R}_{>0}$ is a monotone gain function
(e.g.\ $f(\rho) = e^{g\rho}$ with gain parameter $g \geq 0$). A higher
$\prc_{ij}$ means that agent $i$ weights agent $j$'s actions more heavily when
updating its own $\C_i$ --- precisely the attentional inheritance mechanism
formalised in Section~\ref{sec:installation}.

\paragraph{Key extension over \citet{albarracin2022EpistemicCommunities}.}
Albarracin et al.\ treat all inter-agent channels with a uniform precision;
the community's attractor structure follows from the topology of which sources
agents prefer to sample. Here we add a \emph{prestige layer} that modulates
channel precision independently of source preference. The prestige cue
$\rho_{ij}$ can be operationalised in two ways that correspond to distinct
cultural-learning regimes identified by \citet{henrichGilWhite2001Prestige}:

\begin{enumerate}[label=(\alph*)]
  \item \textbf{Exogenous prestige.} A small fraction $p_h \approx 0.15$ of
    agents are designated high-prestige at the outset ($\rho_{ij} = 1$ for
    high-prestige $j$, $\rho_{ij} = 0$ otherwise). This recovers the simplest
    case of model-based learning studied by \citet{chudek2012BystanderGaze},
    where bystanders differentially attend to marked individuals.

  \item \textbf{Endogenous prestige.} $\rho_{ij}$ self-organises via a
    Henrich-style success-bias rule: after each round, $\rho_{ij}$ is
    proportional to the reduction in free energy that agent $j$ achieved in the
    preceding episode.  Agents that minimise their prediction error most
    efficiently accrue prestige and attract greater attentional weight,
    implementing the ``copy the successful'' heuristic at the precision level.
\end{enumerate}

Both variants share the functional form~\eqref{eq:prestige-precision}; they
differ only in how $\rho_{ij}$ evolves. Variant (b) is the richer case for
the norm-formation question but variant (a) is analytically cleaner for
isolating the regime transitions.

\subsection*{Figure 4 --- Simulation schematic}

\begin{figure}[ht]
\centering
% NOTE: This TikZ figure uses the `graphs` and `graphdrawing` libraries
% (loaded in main.tex) together with the `force` graph-drawing algorithm.
% These require LuaLaTeX. When compiling with pdfLaTeX, replace the
% \graph{} block with manually positioned \node / \draw commands.
\begin{tikzpicture}[
  agent/.style   = {circle, draw, minimum size=0.75cm, inner sep=1pt,
                    font=\small},
  prestige/.style= {circle, draw, very thick, minimum size=0.75cm, inner sep=1pt,
                    font=\small, fill=gray!20},
  chan/.style    = {->, >=Stealth, gray},
  hichan/.style  = {->, >=Stealth, line width=2.5pt, black!70},
  label distance=2pt
]
% --- manually positioned layout (pdfLaTeX-safe version) ---
% Outer ring: 8 ordinary agents
\foreach \i/\angle in {1/0, 2/45, 3/90, 4/135, 5/180, 6/225, 7/270, 8/315}{
  \node[agent] (a\i) at ({2.4*cos(\angle)},{2.4*sin(\angle)})
        {$\C_{\i}$};
}
% Inner core: 2 high-prestige agents
\node[prestige] (p1) at (-0.7, 0.0) {$\C^{*}_{1}$};
\node[prestige] (p2) at ( 0.7, 0.0) {$\C^{*}_{2}$};

% Low-precision channels (ordinary → ordinary, thin)
\foreach \s/\t in {1/2, 2/3, 3/4, 4/5, 5/6, 6/7, 7/8, 8/1, 1/5, 3/7}{
  \draw[chan] (a\s) -- (a\t);
}
% High-precision channels (ordinary → prestige, thick)
\foreach \s in {1,2,3,4,5,6,7,8}{
  \draw[hichan] (a\s) -- (p1);
  \draw[hichan] (a\s) -- (p2);
}
% Mutual channel between high-prestige agents
\draw[hichan, <->] (p1) -- (p2);

% Prestige labels on thick edges (one representative)
\draw[hichan] (a1) -- node[above right, font=\scriptsize]{$\prc_{ij}\!=\!\prc_0 e^{g\rho_{ij}}$} (p1);

% Legend
\node[agent,  below left=0.3cm and 0.0cm of a8, font=\scriptsize]  (leg1) {};
\node[right=0.05cm of leg1, font=\scriptsize] {ordinary agent ($\C_i$)};
\node[prestige, below=0.45cm of leg1, font=\scriptsize] (leg2) {};
\node[right=0.05cm of leg2, font=\scriptsize] {high-prestige agent ($\C^*$)};
\draw[chan,  below=0.6cm of leg2] (0,-3.6) -- node[right, font=\scriptsize]{low $\prc_{ij}$} ++(0.8,0);
\draw[hichan] (0,-3.95) -- node[right, font=\scriptsize]{high $\prc_{ij}$} ++(0.8,0);
\end{tikzpicture}
\caption{Proposed simulation schematic: $N$ agents with individual prior
preferences $\C_i$ observe each other through precision-weighted channels
modulated by prestige cues $\rho_{ij}$.  High-prestige agents (shaded, inner
ring) attract thick observation channels from all other agents; ordinary
agents are connected by thin channels with baseline precision $\prc_0$.  The
community-level $\C$-distribution evolves via attention-mediated copying as
defined in Eq.~\eqref{eq:prestige-precision}.}
\label{fig:sim-schematic}
\end{figure}

\subsection*{Predicted regime transitions}

As the prestige-precision gain $g$ increases from zero, the model predicts
the following sequence of community-level behaviours. At \emph{low} $g$
($g \approx 0$), inter-agent observation channels carry negligible precision
above the individual's self-model; transmitted information is too weak to
overcome agents' prior heterogeneity, and the population $\C$-distribution
remains broad---no norm forms. At \emph{intermediate} $g$, transmission is
strong enough for iterative updating to draw the community toward a single
high-density region: a shared $\C$ emerges, corresponding to the
\emph{shaped} $\C$ regime of Section~\ref{sec:shape-of-C}. At \emph{high} $g$ the
dynamics bifurcate: sub-communities that share initial $\C$-proximity lock
onto distinct attractors, recovering the echo-chamber result of
\citet{albarracin2022EpistemicCommunities} by a complementary mechanism (there,
echo chambers follow from source-sampling bias; here, they follow from
heterogeneous prestige cue weights). At \emph{very high} $g$, the community
collapses onto a single high-prestige agent's $\C$, producing the
\emph{degenerate fixation} regime: the spiky $\C$ of a community captured by
one model.  These four qualitative outcomes span the spiky/broad/shaped $\C$
taxonomy introduced in Section~\ref{sec:shape-of-C} and predicted by the stability
analysis of Section~\ref{sec:loop}.

\subsection*{Quantities to measure}

The simulation collects three primary observables as functions of $g$:

\begin{enumerate}[label=(\roman*)]
  \item \textbf{Mean pairwise KL divergence} $\bar{D}(g) = \binom{N}{2}^{-1}
    \sum_{i < j} \KL{\C_i}{\C_j}$, measuring aggregate heterogeneity of the
    community's prior-preference distribution; this should be non-monotone in
    $g$, peaking at the echo-chamber bifurcation before collapsing at
    degenerate fixation.

  \item \textbf{Population $\C$-entropy} $H(g) = -\int p(\C) \log p(\C)\,
    \mathrm{d}\C$, estimated by kernel density over the empirical distribution
    of $\C_i$ vectors; a unimodal drop in $H$ marks the consensus regime,
    while a bimodal recovery marks the bifurcation.

  \item \textbf{Attractor count} $K(g)$, estimated by $k$-means clustering on
    the converged $\{\C_i\}$ with silhouette-based model selection; $K$
    should trace the sequence $K > 1 \to K = 1 \to K = 2 \to K = 1$ as $g$
    increases through the four regimes.
\end{enumerate}

These observables admit a complementary spectral-graph reading that
connects the simulation directly to the structural analysis of
\citet{heins2024CollectiveSurprise}. Construct the agent influence graph
$\mathcal{G}$ with edge weights equal to the realised $\prc_{ij}$ at
convergence. The \emph{spectral gap} $\lambda_2(\mathcal{G})$ --- the second
eigenvalue of the normalised Laplacian --- measures how quickly information
diffuses across the weakest bottleneck; a large spectral gap indicates a
well-mixed community (consensus regime), while a small spectral gap flags a
bottleneck consistent with bifurcated echo chambers. The \emph{Fiedler
vector} $\mathbf{v}_2(\mathcal{G})$ identifies the natural partition of the
community: its sign pattern should recover the two attractor clusters at high
$g$. Tracking $\lambda_2$ and $\mathbf{v}_2$ as functions of $g$ thus gives
a spectral read-out of the same regime transitions measured by
$\bar{D}(g)$ and $K(g)$.

\subsection*{Connection to the silly-rules result}

The prediction in regime (ii)---the intermediate-$g$ fragile-consensus
window---has a direct analogue in \citet{koster2022SillyRules}. Silly rules,
in the cooperative-AI setting, are arbitrary behavioural taboos that carry no
direct welfare content but substantially improve agents' capacity to learn
stable enforcement and compliance. In the present framework the mechanism is
transparent: silly rules pump precision on the policy prior
``follow-norms-in-general'' without committing to any specific $\C$-content.
They are, in effect, a device for raising $g$ artificially at the level of
policy selection, making the normative channel more legible without specifying
what travels through it. Our framework predicts that silly rules should be
most beneficial precisely in the intermediate-$g$ regime, where the
consensus-forming process is fragile: the channel precision is high enough to
synchronise agents but the attractor is not yet locked in, making additional
precision-pumping most valuable. At very high $g$ (degenerate fixation) silly
rules add noise around an already-collapsed distribution; at low $g$ they
cannot compensate for an insufficiently precise transmission channel. The
intermediate window is where content-neutral precision amplification does
meaningful normative work, consistent with the \citeauthor{koster2022SillyRules}
finding that spurious normativity enhances acquisition of welfare-relevant
compliance.



%% ============================================================
\section{Discussion and Open Work}
\label{sec:discussion}
%% ============================================================

We have argued that cooperative AI's working concept of norm is best refined
as a stabilised region of high-precision shared prior preference, installed
in agents by a precision-weighted attentional channel that prestige cues and
credibility-enhancing displays modulate. The fixed-point loop we formalise
--- community priors shape salience, salience directs copying, copying
installs priors --- is the same loop that the cultural-evolution tradition
documents informally and that active inference has the apparatus to render
formally; the contribution is to close the loop explicitly and to locate it
on the prior-preference content layer ($\C$), one level up from the
policy-precision account of social power developed by
\citet{albarracin2025SocialPower}.

Three threads warrant brief comment.

\paragraph{Gradual disempowerment as channel capture.}
The substrate of normative life, on this account, is the precision channel
through which $\C$ is installed; whoever mediates that channel shapes which
priors get installed in the next generation. The
\citet{kulveit2025GradualDisempowerment} thesis acquires a clean mechanistic
reading in this picture: the worry is not that AI systems will explicitly
override human preferences but that they already mediate the channel through
which preferences are formed. Recommendation systems, conversational
interfaces, and AI-shaped attention markets are precision-allocators, not
just content-providers. Whether their integration constitutes capture,
augmentation, or genuine partnership turns on whether the
$\C$-installation channel remains a function of the community's own values
or becomes a function of the system's optimisation objective. We flag this
as a first-order question for cooperative AI without developing it here.

\paragraph{Limitations and what is not in this paper.}
Several formal moves require empirical anchoring we do not provide. The
CRED-as-precision-modulator hypothesis is novel and predicts specific
dose-response relationships between display costliness and posterior
precision on transmitted values; the simulation in
Section~\ref{sec:simulation} tests one consequence (regime transitions in
prestige-precision gain) but not this empirical specification directly. A
follow-up empirical programme would profitably target precision-shift
markers under controlled CRED exposure. Separately, this paper does not
develop the structural / spectral reading hinted at in
Section~\ref{sec:simulation}'s in-community measurements: the broader
project of using spectral-graph quantities as cross-mechanism signatures
of coordination structure is left for a separate paper. Finally, the
political-theory bridge (Bourdieu's habitus, Lukes' third face of power as
preference-shaping) is invoked as motivation but not formalised here beyond
the explicit differentiation from \citet{albarracin2025SocialPower}.

\paragraph{What we hope this draft does.}
We have presented this as a working draft for active-inference colleagues.
We welcome feedback, especially from those embedded in multi-agent
active-inference, cultural-evolution, or cooperative-AI norm research,
where the framework most directly intersects existing programmes and where
the formalism is most likely to over- or under-reach.

\bibliographystyle{plainnat}
\bibliography{values-salience}

\end{document}
